\section{Marco Teórico}

En esta sección se presentan los fundamentos técnicos y conceptuales que sustentan el desarrollo del presente proyecto. Se abordan tecnologías clave del sector automotriz, enfocadas principalmente en la comunicación interna entre sistemas electrónicos.

\subsection{Protocolo CAN BUS}

El \textbf{Controller Area Network} (CAN BUS) es un protocolo de comunicación serial creado por Bosch en los años 80, diseñado para aplicaciones en la automoción y estandarizado posteriormente en la norma ISO 11898 \cite{iso11898}. Su característica esencial es permitir que múltiples microcontroladores y dispositivos (nodos) se comuniquen entre sí sin necesidad de un maestro central.\\

La robustez del protocolo se debe a su esquema de arbitraje: cuando dos nodos transmiten al mismo tiempo, prevalece el mensaje con identificador de menor valor, lo que asegura una comunicación determinista y priorizada. Además, utiliza comunicación diferencial mediante dos líneas (CAN\_H y CAN\_L), alcanzando velocidades típicas de hasta 1 Mbps.\\

El bus CAN utiliza un protocolo \textbf{basado en mensajes}, donde los datos se transmiten en tramas en lugar de conexiones punto a punto. 

Cada trama de mensaje CAN consta de varios campos: el \textbf{Inicio de Trama} (SOF), un Identificador (ID), Campo de Control, Campo de Datos, \textbf{Verificación por Redundancia Cíclica} (CRC), Campo de Reconocimiento (ACK) y el Fin de Trama (EOF). El Identificador es crucial, ya que determina la prioridad del mensaje; los valores numéricos más bajos tienen mayor prioridad. Esto permite que los mensajes críticos, como los relacionados con el frenado o el control del motor, se transmitan antes que datos menos urgentes. \\

El protocolo CAN utiliza un mecanismo de Acceso Múltiple por \textbf{Detección de Portadora con Detección de Colisiones} (CSMA/CD), pero con un proceso de arbitraje único. Cuando varios nodos intentan transmitir al mismo tiempo, el nodo con el ID más bajo gana la arbitraje y continúa transmitiendo, mientras que los demás se retiran automáticamente, lo que garantiza que no se pierda ningún dato.\\

\subsubsection{Tramas}
El bus Controller Area Network (CAN) utiliza una estructura específica de tramas de mensajes para la transmisión de datos, lo que garantiza una comunicación eficiente y confiable entre las unidades de control electrónico (ECUs) en entornos como los sistemas automotrices e industriales.
El protocolo CAN define varios tipos de tramas: \textbf{tramas de datos, tramas remotas, tramas de error y tramas de sobrecarga}, cada una con una función específica.\\
\subsubsection*{Trama de datos}
Una trama de datos se utiliza para transmitir información desde un nodo a uno o más nodos receptores y consta de varios campos.
Comienza con el Inicio de Trama (\textbf{SOF}), un solo bit dominante (lógica 0) que señala el comienzo de un nuevo mensaje y sincroniza todos los nodos en el bus.
Después del SOF, sigue el \textbf{Campo de Arbitraje}, que contiene el ID del mensaje (11 bits en formato estándar, permitiendo 2048 identificadores únicos) y el bit \textbf{Remote Transmission Request (RTR)}.
El bit RTR distingue entre una trama de datos (RTR = 0) y una trama remota (RTR = 1).
El ID del mensaje también determina la prioridad del mensaje durante la arbitraje del bus; \textbf{un ID numéricamente más bajo tiene mayor prioridad.}\\

A continuación se encuentra el Campo de Control, que incluye bits reservados (R0, R1) y el \textbf{Código de Longitud de Datos (DLC)}, un campo de 4 bits que indica la cantidad de bytes de datos (de 0 a 8) en la trama.
Le sigue el Campo de Datos, que contiene la carga útil del mensaje, con un máximo de 8 bytes en el CAN estándar.
Después de los datos, un campo de \textbf{Verificación por Redundancia Cíclica (CRC)} de 15 bits garantiza la detección de errores; el receptor calcula su propio CRC y lo compara con el valor transmitido.
Una discrepancia indica un error de transmisión.\\

El \textbf{Campo de Reconocimiento (ACK)} consta de una ranura ACK y un delimitador ACK.
El nodo transmisor envía un bit recesivo en la ranura ACK, que los nodos receptores sobrescriben con un bit dominante (lógica 0) para confirmar la recepción exitosa; si esto no ocurre, se produce un error de ACK.
La trama concluye con el \textbf{ Fin de Trama (EOF)}, siete bits recesivos que marcan el final de la trama, seguidos por un espacio de Intermisión (IFS) de 3 bits, tras lo cual el bus vuelve al estado inactivo y cualquier nodo puede iniciar una nueva transmisión.\\
\subsubsection*{Tramas remotas}
Las tramas remotas son estructuralmente similares a las tramas de datos, pero carecen del campo de datos; se utilizan para que un nodo solicite datos a otro nodo, que luego responde con una trama de datos usando el mismo ID de mensaje.

\subsubsection*{Tramas de error}
Las tramas de error son transmitidas por cualquier nodo que detecta un error y sirven para abortar la trama actual, mientras que las tramas de sobrecarga introducen una demora entre tramas para gestionar la carga de procesamiento de los nodos.\\


\subsection{Unidades de Control Electrónico (ECUs)}

Las ECUs son microcontroladores dedicados al control de funciones específicas dentro del vehículo, como el motor, la transmisión, los frenos, la climatización o el sistema de infoentretenimiento. Cada ECU intercambia información con otras a través del bus CAN, conformando una arquitectura de red interna. \\

En automóviles modernos, el número de ECUs puede superar las 50 unidades, reflejando la creciente complejidad electrónica. En contraste, en vehículos más antiguos las funciones electrónicas eran más aisladas y, en muchos casos, limitadas al diagnóstico del motor.

\subsection{El protocolo ISP (In-System Programming)}
El \textbf{ISP} (\textit{In-System Programming}) es un protocolo de programación que permite cargar firmware en microcontroladores directamente en la placa donde se encuentran montados, sin necesidad de retirarlos o emplear programadores externos especializados. Este enfoque resulta fundamental en el desarrollo de sistemas embebidos, ya que facilita la iteración rápida, la actualización de software y el mantenimiento de los dispositivos una vez desplegados en el entorno de aplicación.

El principio básico de ISP consiste en habilitar al microcontrolador para que, mediante una interfaz de comunicación serial dedicada, reciba instrucciones de programación desde un dispositivo externo (generalmente un computador o un programador USB) y almacene el código binario en su memoria flash. Para ello, el microcontrolador dispone de un \textbf{bootloader} o de rutinas internas específicas que gestionan este proceso.

\subsubsection*{Características principales}

\begin{itemize}
    \item \textbf{Acceso directo a la memoria:} permite escribir y reescribir el contenido de la memoria de programa sin desmontar el microcontrolador.
    \item \textbf{Uso de pines estándar:} típicamente se emplean líneas como MOSI, MISO, SCK y RESET, compartidas con la interfaz SPI del dispositivo.
    \item \textbf{Compatibilidad amplia:} la mayoría de los microcontroladores de las familias \textbf{Atmel AVR}, \textbf{PIC} y algunos ARM Cortex incluyen soporte para ISP.
    \item \textbf{Velocidad de programación:} el proceso es relativamente rápido debido a la naturaleza serial síncrona de la comunicación.
\end{itemize}

\subsubsection*{Aplicaciones en el proyecto}

En el contexto del presente trabajo, el protocolo ISP será utilizado para programar los microcontroladores \textbf{Atmega} que actúan como nodos en la red CAN. Estos microcontroladores, equipados con el chip MCP2515 para conectividad, requieren que su firmware se cargue y actualice durante el ciclo de desarrollo y pruebas. La posibilidad de reprogramarlos en la propia tarjeta sin desoldar ni extraer componentes representa una ventaja significativa en términos de agilidad y reducción de costos.

Asimismo, el uso de ISP permite realizar ajustes incrementales en los algoritmos de interpretación de mensajes CAN, control de sensores o activación de actuadores, garantizando un flujo de desarrollo iterativo y flexible. Esta capacidad resulta clave en proyectos de I+D donde la experimentación y la mejora continua son parte inherente del proceso.

\subsection{Protocolo SPI (Serial Peripheral Interface)}

El \textbf{SPI} (\textit{Serial Peripheral Interface}) es un protocolo de comunicación serial síncrono desarrollado por Motorola a mediados de los años 80, ampliamente utilizado en sistemas embebidos para la interconexión de microcontroladores con periféricos de alta velocidad, tales como memorias, sensores, convertidores A/D y controladores de comunicación como el \textbf{MCP2515} utilizado en este proyecto. Su diseño prioriza la \textbf{velocidad y simplicidad}, lo que lo convierte en una de las interfaces más empleadas en aplicaciones que requieren intercambio rápido y confiable de datos.

\subsubsection*{Principios de funcionamiento}

El protocolo SPI se basa en un esquema \textbf{maestro–esclavo}, donde un dispositivo maestro controla el flujo de comunicación mediante una señal de reloj compartida. Para ello, utiliza cuatro líneas principales:

\begin{itemize}
    \item \textbf{MOSI (Master Out, Slave In)}: línea de datos desde el maestro hacia el esclavo.
    \item \textbf{MISO (Master In, Slave Out)}: línea de datos desde el esclavo hacia el maestro.
    \item \textbf{SCK (Serial Clock)}: señal de reloj generada por el maestro para sincronizar la transferencia.
    \item \textbf{SS (Slave Select)}: señal para habilitar o deshabilitar un esclavo específico en el bus.
\end{itemize}

La transferencia de datos se realiza de forma \textbf{full-duplex}, lo que significa que mientras el maestro envía un bit por MOSI, simultáneamente recibe otro desde MISO. Esto permite velocidades de comunicación superiores a las de protocolos como I\textsuperscript{2}C o UART.

\subsubsection*{Características principales}

\begin{itemize}
    \item \textbf{Alta velocidad de transferencia:} puede alcanzar decenas de MHz en implementaciones típicas.
    \item \textbf{Simplicidad de protocolo:} no requiere direccionamiento complejo ni tramas extensas.
    \item \textbf{Full-duplex:} permite transmisión y recepción simultánea de datos.
    \item \textbf{Flexibilidad en configuración:} parámetros como polaridad y fase del reloj (CPOL y CPHA) pueden ajustarse para adaptarse a distintos periféricos.
    \item \textbf{Limitación en número de esclavos:} cada dispositivo requiere su propia línea de selección (SS), lo que dificulta la escalabilidad en sistemas con muchos periféricos.
\end{itemize}

\subsubsection*{Aplicaciones en el proyecto}

En el presente proyecto, el protocolo SPI se emplea como interfaz de comunicación entre los microcontroladores \textbf{Atmega} y el controlador CAN \textbf{MCP2515}. Gracias a la velocidad de SPI, es posible enviar y recibir tramas CAN de forma eficiente, asegurando que los nodos diseñados puedan integrarse correctamente en la red del vehículo.

Además, el uso de SPI permite que los microcontroladores intercambien información con otros periféricos de apoyo, como memorias externas o sensores de alta velocidad, en caso de ser necesario. De esta forma, SPI actúa como un \textbf{puente fundamental} entre la lógica de los microcontroladores y el hardware de comunicación que conecta al vehículo con la red CAN distribuida.

\subsection{Protocolo I\textsuperscript{2}C (Inter-Integrated Circuit)}
El protocolo \textbf{I\textsuperscript{2}C} (\textit{Inter-Integrated Circuit}) es un estándar de comunicación serial desarrollado por Philips a principios de los años 80, diseñado para permitir la interconexión sencilla de múltiples dispositivos electrónicos empleando únicamente dos líneas de comunicación: \textbf{SDA} (Serial Data) y \textbf{SCL} (Serial Clock). Su filosofía se centra en la simplicidad de conexión, bajo consumo de pines y capacidad de escalar a configuraciones maestro-esclavo en sistemas embebidos.

\subsubsection*{Principios de funcionamiento}

El bus I\textsuperscript{2}C permite que un \textbf{dispositivo maestro} controle la comunicación, generando las señales de reloj (SCL) y enviando las instrucciones a uno o más \textbf{dispositivos esclavos}. Cada esclavo conectado al bus dispone de una dirección única, lo que hace posible que múltiples circuitos integrados (sensores, memorias, conversores A/D, etc.) compartan las mismas líneas físicas sin interferencia. 

El protocolo es síncrono y soporta velocidades de transferencia que van desde:
\begin{itemize}
    \item 100 kHz (modo estándar),
    \item 400 kHz (modo rápido),
    \item hasta varios MHz en implementaciones extendidas (Fast-mode Plus, High-speed mode).
\end{itemize}

\subsubsection*{Características principales}

\begin{itemize}
    \item \textbf{Uso eficiente de pines:} solo requiere dos líneas para interconectar múltiples dispositivos.
    \item \textbf{Escalabilidad:} admite múltiples esclavos y, en algunas implementaciones, múltiples maestros.
    \item \textbf{Direccionamiento:} cada esclavo posee una dirección de 7 o 10 bits, lo que posibilita la conexión de decenas de periféricos en un mismo bus.
    \item \textbf{Comunicación half-duplex:} en la misma línea SDA se transmiten y reciben datos, bajo control del maestro.
    \item \textbf{Soporte en hardware:} gran variedad de microcontroladores incluyen módulos I\textsuperscript{2}C dedicados.
\end{itemize}

\subsubsection*{Aplicaciones en el proyecto}

En el presente proyecto, el protocolo I\textsuperscript{2}C se considera clave para la \textbf{interconexión local de sensores} dentro de cada nodo controlado por microcontroladores Atmega o Raspberry Pi. Por ejemplo, los sensores de \textbf{temperatura, humedad y presión ambiental} suelen disponer de interfaces I\textsuperscript{2}C, lo que simplifica su integración al requerir solo dos pines del microcontrolador.

De este modo, mientras que el \textbf{CAN BUS} sirve como red troncal de comunicación entre las distintas ECUs del vehículo, el \textbf{I\textsuperscript{2}C} funciona como un sub-bus interno dentro de cada ECU, facilitando la conexión de múltiples periféricos sin necesidad de recurrir a buses adicionales más complejos. Esta complementariedad entre ambos protocolos asegura un diseño modular y escalable del sistema.
\subsection{Protocolo USART (Universal Synchronous/Asynchronous Receiver-Transmitter)}

El \textbf{USART} (\textit{Universal Synchronous/Asynchronous Receiver-Transmitter}) es un módulo de comunicación serial integrado en numerosos microcontroladores, cuya función es transmitir y recibir datos de manera eficiente mediante una interfaz configurable tanto en modo síncrono como en modo asíncrono. A diferencia de buses como I\textsuperscript{2}C o SPI, el USART no define un protocolo completo de interconexión, sino que proporciona una \textbf{interfaz física y lógica} para la transmisión de datos serializados entre dispositivos.

\subsubsection*{Principios de funcionamiento}

En su modo más común, el \textbf{asíncrono}, el USART implementa la comunicación serial UART tradicional, donde los datos se transmiten bit a bit a través de dos líneas:
\begin{itemize}
    \item \textbf{TX (Transmit)}: línea de transmisión de datos.
    \item \textbf{RX (Receive)}: línea de recepción de datos.
\end{itemize}

El flujo de datos está organizado en \textbf{tramas} que incluyen bits de inicio, bits de datos, opcionalmente un bit de paridad y uno o más bits de parada. Esta estructura permite que el receptor sincronice la comunicación sin necesidad de compartir una señal de reloj, aunque ambas partes deben configurarse a la misma velocidad de transmisión (\textit{baud rate}).

En el modo \textbf{síncrono}, el USART emplea además una línea de reloj compartida que asegura la sincronización exacta entre transmisor y receptor, lo que permite mayores velocidades de transferencia y una comunicación más robusta.

\subsubsection*{Características principales}

\begin{itemize}
    \item \textbf{Flexibilidad:} soporta tanto transmisión síncrona como asíncrona.
    \item \textbf{Velocidad:} puede operar en un amplio rango de tasas de transmisión, típicamente entre 300 baudios y varios Mbps.
    \item \textbf{Eficiencia:} requiere solo dos líneas en modo asíncrono (TX y RX), y tres en modo síncrono (TX, RX y CLK).
    \item \textbf{Compatibilidad:} base de estándares de comunicación ampliamente usados como RS-232, RS-485 y TTL serial.
    \item \textbf{Implementación simple:} disponible de forma nativa en la mayoría de microcontroladores.
\end{itemize}

\subsubsection*{Aplicaciones en el proyecto}

En este proyecto, el uso del protocolo USART se contempla principalmente para:
\begin{itemize}
    \item \textbf{Depuración y pruebas}: permitir la comunicación entre los microcontroladores (Atmega o Raspberry Pi) y un computador externo para monitorear mensajes, estados del bus CAN y parámetros de sensores.
    \item \textbf{Configuración inicial}: facilitar la carga de parámetros de red o ajustes de firmware durante la fase de pruebas en banco.
    \item \textbf{Interconexión de módulos externos}: algunos periféricos que no soportan I\textsuperscript{2}C pueden integrarse mediante interfaces UART, garantizando versatilidad en la expansión del sistema.
\end{itemize}

Gracias a su simplicidad y soporte extendido, el USART constituye una herramienta esencial en el ciclo de desarrollo, complementando al bus CAN (para comunicación distribuida en el vehículo) y a I\textsuperscript{2}C (para la integración de sensores locales en cada nodo).


\subsection{El MCP2515 como controlador CAN}

Para integrar nodos CAN en sistemas embebidos que no cuentan con este protocolo de manera nativa, se emplean controladores externos como el \textbf{MCP2515}, fabricado por Microchip Technology. Este chip, conectado mediante interfaz SPI, implementa el protocolo CAN 2.0B y facilita la transmisión y recepción de mensajes.\\

Entre sus características destacan la capacidad de trabajar a 1 Mbps, tres buffers de transmisión y dos de recepción, así como filtros y máscaras programables que permiten gestionar eficientemente el tráfico en el bus. En este proyecto, el MCP2515 actúa como intermediario entre los microcontroladores y el bus CAN, garantizando escalabilidad y compatibilidad con estándares de la industria.

\subsection{Microcontroladores y plataformas embebidas utilizadas}

El proyecto contempla el uso de microcontroladores \textbf{Atmega} (como los empleados en placas Arduino) para funciones de control en tiempo real y bajo consumo, en conjunto con plataformas más potentes como \textbf{Raspberry Pi} para la gestión multimedia y de conectividad.\\

Los Atmega destacan por su sencillez y amplia documentación, mientras que las Raspberry Pi permiten manejar interfaces gráficas, procesamiento de audio y comunicación con redes externas. Esta combinación asegura una arquitectura balanceada entre nodos ligeros y nodos avanzados.

\subsection{Arquitectura distribuida en automoción}

La arquitectura distribuida en vehículos implica que cada función relevante es gestionada por una ECU especializada, todas interconectadas por un bus común. Este enfoque modular mejora la escalabilidad, la tolerancia a fallos y la posibilidad de añadir nuevas funciones sin rediseñar completamente el sistema.\\

Nuestro proyecto se inspira en este principio para añadir dos nuevas unidades: una de confort y otra de multimedia, que se integrarán al bus CAN extendiendo las capacidades del vehículo base.

\subsection{Diagnóstico mediante OBD-II}

El \textbf{On-Board Diagnostics II} (OBD-II) es un estándar internacional (SAE J1962) que permite acceder a datos de diagnóstico del motor y emisiones. Aunque puede emplear diferentes protocolos (ISO 9141, KWP2000, J1850), desde 2008 la normativa exige que todos los vehículos ligeros vendidos en Estados Unidos soporten comunicación CAN en este conector.\\

En el caso del vehículo de estudio, la única comunicación CAN existente se limita al OBD-II para diagnóstico. Por ello, uno de los objetivos principales del proyecto es extender la red interna para habilitar un bus CAN activo y bidireccional, capaz de integrar nuevas ECUs con funcionalidades ampliadas.
